\begin{corollary}[Coarse span consequence]\label{cor:energy}
For every \(m\ge7\) and \(y_1<\cdots<y_m\),
\[
 E_m+\beta(y_m-y_1)\ge\delta(m-6).
\]
For the certified vector \eqref{eq:p},
\[
 E_m+\frac1{500}(y_m-y_1)
 \ge\frac{39}{10000}(m-6).
\]
\end{corollary}

\begin{proof}
Since \(q_r\le\beta\),
\[
 P_B=\sum_rq_rg_r
 \le\beta\sum_rg_r
 =\beta(y_m-y_1).
\]
Apply Proposition~\ref{prop:redist}.
\end{proof}


\section{Pressure-preserving block stability}\label{sec:blockstability}

Put
\[
 \delta:=\frac{39}{10000},
 \qquad
 \beta:=\frac1{500},
 \qquad
 A_m:=\delta(m-6),
 \qquad
 Q_m:=\beta(m-6).
\]
For a block \(B=(y_1,\ldots,y_m)\), let \(P_B\) be the exact pressure
\eqref{eq:qr}.  Proposition~\ref{prop:redist} gives
\begin{equation}\label{eq:EplusP}
 E_m+P_B\ge A_m,
 \qquad
 \sum_{r=1}^{m-1}q_r=Q_m.
\end{equation}

\begin{remark}[Immediate accounting gain]\label{rem:accounting-gain}
Take \(m=262\).  Then
\[
 A_{262}=\frac{624}{625}<1,
 \qquad
 Q_{262}=\frac{64}{125}.
\]
The proof of the old block-stability lemma applies verbatim with \(P_B\) in
place of \(\beta\operatorname{span}(B)\): if \(P_B\ge A_{262}\) the claim is
trivial, while if \(P_B<A_{262}\), the positive lower bound
\(q_r\ge\min_jp_j\) confines the block span to a fixed compact interval and
Lemma~\ref{lem:block} together with \eqref{eq:EplusP} gives
\[
 \mathcal D(G_B)+P_B\ge A_{262}-o(1).
\]
Averaging the exact pressure as in Section~\ref{sec:pinching} then yields
\[
 \liminf_{T\to\infty}\frac{N_0^s(T,2T)}{N(T,2T)}
 \ge
 \frac{655000H_{\rm MT}-1280}{652504}
 =
 0.6731115225500757511923586\ldots .
\]
Thus pressure preservation alone already improves the previous theorem.
\end{remark}


\subsection{Adjacent-pair pinching}

Let \(K=(k(y_i-y_j))_{1\le i,j\le m}\) be a limiting Gram matrix.  Then
\(K\succeq0\), \(K_{ii}=1\), and
\begin{equation}\label{eq:Em-trace}
 E_m
 =
 \operatorname{tr}(K-I)^2
 =
 2\sum_{1\le i<j\le m}w(y_j-y_i).
\end{equation}

\begin{lemma}[Adjacent-pair defect]\label{lem:adjacent-pairs}
For consecutive gaps \(g_j=y_{j+1}-y_j\),
\begin{equation}\label{eq:adjacent-pair-defect}
 \boxed{\mathcal D(K)\ge\sum_{j=1}^{m-1}w(g_j)}.
\end{equation}
\end{lemma}

\begin{proof}
Pinch \(K\) first into the principal \(2\times2\) blocks
\[
 (1,2),(3,4),\ldots
\]
and then into
\[
 (2,3),(4,5),\ldots .
\]
The spectral trace functional \(X\mapsto\operatorname{tr}\Psi(X)\) cannot
increase under either pinching, since pinching is a random-unitary average
and this trace functional is convex and unitarily invariant.

The block associated with a consecutive gap \(g\) is
\[
 \begin{pmatrix}1&k(g)\\ \overline{k(g)}&1\end{pmatrix},
\]
whose eigenvalues are \(1\pm|k(g)|\in[0,2]\).  Its defect is therefore
exactly \(2|k(g)|^2=2w(g)\).  The first pinching gives
\[
 \mathcal D(K)\ge2\sum_{j\ {\rm odd}}w(g_j),
\]
and the second gives the corresponding inequality over even \(j\).  Averaging
the two inequalities proves \eqref{eq:adjacent-pair-defect}.
\end{proof}


\subsection{A scalar kernel inequality}

Set
\begin{equation}\label{eq:g0t}
 g_0:=\frac{89}{100},
 \qquad
 t:=\frac{33}{2}.
\end{equation}

\begin{lemma}[Scalar pressure bound]\label{lem:scalar-pressure}
For every \(g\ge0\) and every \(q\in[0,\beta]\),
\begin{equation}\label{eq:scalar-pressure}
 w(g)+tqg\ge tqg_0.
\end{equation}
\end{lemma}

\begin{proof}
The normalized kernel has the integral representation
\[
 k(g)
 =
 \int_{-1/2}^{1/2}
 \frac{\cos(\sqrt2\,u)}
      {\sqrt2\sin(1/\sqrt2)}
 \cos(2\pi gu)\,du .
\]
For \(0<g\le1\), differentiation under the integral sign gives
\(k'(g)<0\): on \(0<u<1/2\), the density is positive and
\(\sin(2\pi gu)>0\).  Moreover
\[
 k(1)=\frac1{2\pi^2-1}>0.
\]
Hence \(w(g)=k(g)^2\) is decreasing on \([0,1]\).

The exact-arithmetic verifier accompanying this version proves
\begin{equation}\label{eq:w089}
 w\!\left(\frac{89}{100}\right)
 >
 \frac{2937}{100000}
 =
 t\beta g_0.
\end{equation}
If \(0\le g\le g_0\), then
\[
 w(g)\ge w(g_0)>t\beta g_0\ge tqg_0.
\]
If \(g\ge g_0\), then \(tqg\ge tqg_0\).  This proves
\eqref{eq:scalar-pressure}.
\end{proof}


\begin{lemma}[Spectral threshold]\label{lem:spectral-threshold}
Let \(K\succeq0\) be \(m\times m\) with diagonal \(1\), and set
\[
 E:=\operatorname{tr}(K-I)^2.
\]
If \(\mathcal D(K)<E\), then
\begin{equation}\label{eq:spectral-threshold}
 \mathcal D(K)>\frac{m}{m-1}.
\end{equation}
\end{lemma}

\begin{proof}
If all eigenvalues of \(K\) belonged to \([0,2]\), then
\(\Psi(\lambda)=(\lambda-1)^2\) for every eigenvalue and
\(\mathcal D(K)=E\).  Hence \(\mathcal D(K)<E\) implies the existence of
an eigenvalue \(1+a>2\), with \(a>1\).

The remaining \(m-1\) eigenvalues have sum \(m-1-a\), hence average
\(1-a/(m-1)\).  Since \(\Psi\) is convex,
\[
 \mathcal D(K)
 \ge
 2a-1
 +(m-1)\Psi\!\left(1-\frac{a}{m-1}\right)
 =
 2a-1+\frac{a^2}{m-1}.
\]
At \(a=1\) the last expression equals \(m/(m-1)\), and it is strictly
increasing for \(a>1\), proving \eqref{eq:spectral-threshold}.
\end{proof}


\begin{proposition}[Strong block stability at \(m=450\)]
\label{prop:blockstab450}
Let \(B\) be a block of
\[
 m=450
\]
consecutive retained simple zeros, with exact pressure \(P_B\) from
\eqref{eq:qr}.  Put
\begin{equation}\label{eq:AQ450}
 A:=A_{450}=\frac{4329}{2500},
 \qquad
 Q:=Q_{450}=\frac{111}{125}.
\end{equation}
Uniformly over such blocks,
\begin{equation}\label{eq:blockstab450}
 \boxed{\mathcal D(G_B)+P_B\ge A-o(1)}.
\end{equation}
\end{proposition}

\begin{proof}
If \(P_B\ge A\), the assertion follows from \(\mathcal D(G_B)\ge0\).
Assume \(P_B<A\).  Since \(q_r\ge\min_jp_j=2714/10^7\),
\[
 \operatorname{span}(B)
 =
 \sum_rg_r
 \le
 \frac{P_B}{\min_jp_j}
 <
 \frac{A}{\min_jp_j}.
\]
Thus all pair separations remain in a fixed compact interval.  Since \(m\) is
fixed, Lemma~\ref{lem:kernel} gives, uniformly in this branch,
\[
 G_B=K_B+o(1)
\]
in matrix norm, where
\[
 K_B=(k(y_i-y_j))_{i,j=1}^m.
\]
By continuity of the spectral trace functional,
\[
 \mathcal D(G_B)=\mathcal D(K_B)+o(1).
\]
It remains to prove
\[
 \mathcal D(K_B)+P_B\ge A.
\]

Write \(K=K_B\) and \(D=\mathcal D(K)\).  By
Proposition~\ref{prop:redist} and \eqref{eq:Em-trace},
\[
 E+P_B\ge A,
 \qquad
 E=\operatorname{tr}(K-I)^2.
\]
Lemma~\ref{lem:adjacent-pairs} and
Lemma~\ref{lem:scalar-pressure} give
\begin{equation}\label{eq:Dplus-tP-lower}
 D+tP_B
 \ge
 \sum_{j=1}^{m-1}\bigl(w(g_j)+tq_jg_j\bigr)
 \ge
 tg_0\sum_{j=1}^{m-1}q_j
 =
 tg_0Q.
\end{equation}

Suppose for contradiction that \(D+P_B<A\).  Since \(E+P_B\ge A\),
we have \(D<E\); Lemma~\ref{lem:spectral-threshold} therefore yields
\[
 D>\frac{m}{m-1}.
\]
Consequently
\begin{align}
 D+tP_B
 &=
 t(D+P_B)-(t-1)D \nonumber\\
 &<
 tA-(t-1)\frac{m}{m-1}.
 \label{eq:Dplus-tP-upper}
\end{align}
For \(m=450\), the lower bound
\eqref{eq:Dplus-tP-lower} contradicts
\eqref{eq:Dplus-tP-upper}, because the exact rational margin is
\begin{equation}\label{eq:contradiction-margin}
 g_0Q
 +
 \left(1-\frac1t\right)\frac{450}{449}
 -A
 =
 \frac{6363}{30868750}
 >0.
\end{equation}
Hence \(D+P_B\ge A\), and the uniform finite-\(T\) statement
\eqref{eq:blockstab450} follows.
\end{proof}


\section{Shifted block pinching with exact pressure}\label{sec:pinching}

Let
\[
 x_1<\cdots<x_{S^\circ}
\]
be the normalized ordinates of the retained central set
\(\mathcal S^\circ\).  As before,
\[
 S^\circ=N_0^s(T,2T)-o(N(T,2T)).
\]

For each of the \(m=450\) offsets modulo \(m\), partition the ordered points
into full consecutive blocks of size \(m\), with only \(O(m)\) endpoint
points left over.  Pinching \(M^\circ\) to the full principal blocks and the
remainder gives
\[
 \mathcal D(M^\circ)\ge\sum_B\mathcal D(G_B).
\]
Summing over all offsets and applying Proposition~\ref{prop:blockstab450},
\begin{equation}\label{eq:shifted-exact}
 m\mathcal D(M^\circ)
 \ge
 (S^\circ-m+1)A
 -\sum_BP_B
 -o(N).
\end{equation}

The pressure term retains its exact positional mass.  For each fixed relative
position \(r\), a given global gap can occupy that position in at most one
sliding \(m\)-point block.  Therefore its coefficient in the sum over all
sliding blocks is at most
\[
 \sum_{r=1}^{m-1}q_r=Q.
\]
Hence
\begin{equation}\label{eq:pressure-global}
 \sum_BP_B
 \le
 Q(x_{S^\circ}-x_1)
 \le
 QN(T,2T)+o(N(T,2T)),
\end{equation}
where the last inequality is the same normalized-span consequence of the
Riemann--von Mangoldt formula used previously.

Combining \eqref{eq:shifted-exact} and \eqref{eq:pressure-global},
and absorbing \(O(1)\) endpoint terms,
\begin{equation}\label{eq:defectglobal-new}
 \mathcal D(M^\circ)
 \ge
 \frac{A}{m}N_0^s(T,2T)
 -\frac{Q}{m}N(T,2T)
 -o(N(T,2T)).
\end{equation}
For \(m=450\),
\begin{equation}\label{eq:defectnumbers-new}
 \mathcal D(M^\circ)
 \ge
 \frac{481}{125000}N_0^s(T,2T)
 -\frac{37}{18750}N(T,2T)
 -o(N(T,2T)).
\end{equation}


\begin{proof}[Proof of Theorem~\ref{thm:main}]
Write
\[
 S=N_0^s(T,2T),
 \qquad
 N=N(T,2T).
\]
Insert \eqref{eq:defectnumbers-new} into
Corollary~\ref{cor:global}:
\[
 S
 \ge
 H_{\rm MT}N
 +\frac{481}{125000}S
 -\frac{37}{18750}N
 -o(N).
\]
Thus
\[
 \frac{124519}{125000}S
 \ge
 \left(H_{\rm MT}-\frac{37}{18750}\right)N-o(N).
\]
Dividing by \(N\) and passing to the lower limit gives
\begin{align*}
 \liminf_{T\to\infty}\frac{S}{N}
 &\ge
 \frac{450H_{\rm MT}-111/125}
      {450-4329/2500}\\
 &=
 \frac{1125000H_{\rm MT}-2220}{1120671}\\
 &=
 0.6731175265883904388095857434106058666\ldots .
\end{align*}
\end{proof}
